\documentclass[10pt,a4paper]{amsart}
\usepackage[margin=19mm]{geometry}
\usepackage{amsmath,amssymb,mathtools}
\usepackage[scr=rsfs,cal=euler]{mathalfa}
\usepackage{xfrac}
\usepackage[unicode,hidelinks,bookmarks=false]{hyperref}
\newcommand{\ie}{i.e.\ }
\newcommand{\cf}{cf.\ }
\newcommand{\resp}{respectively\ }
\newcommand{\Subsection}{\subsection}
\providecommand{\nobreakdash}{\nobreak\hbox{-}}
\setlength{\emergencystretch}{2em}
\multlinegap0pt
\usepackage{bm}
\usepackage{booktabs}
\usepackage{nicefrac}
\usepackage{mathtools}
\usepackage{bbm}
\usepackage{amssymb}
\usepackage{multirow}
\usepackage[shortlabels]{enumitem}
\usepackage{caption}
\usepackage{braket}
\usepackage{siunitx}
\usepackage{algorithm}\usepackage{algpseudocode}
\usepackage{mdframed}
\usepackage{tcolorbox}
\newtheorem{prop}{Proposition}
\newcommand{\E}{\mathbb{E}}
\newcommand{\MeasuredCISet}{\mathcal{D}_{\feas{}}}
\newcommand{\feas}{\mathfrak{f}}
\title{Solving Parameter-Robust Avoid Problems with Unknown Feasibility using Reinforcement Learning}
\author{Oswin So, Eric Yang Yu, Songyuan Zhang, Matthew Cleaveland, Mitchell Black, Chuchu Fan}
\date{Source: arXiv:2602.15817v1 (CC BY 4.0)}
\begin{document}
\maketitle

\noindent\emph{Source and licence.} Solving Parameter-Robust Avoid Problems with Unknown Feasibility using Reinforcement Learning. Version arXiv:2602.15817v1 (CC BY 4.0), distributed by arXiv under the Creative Commons Attribution 4.0 International licence, http://creativecommons.org/licenses/by/4.0/, as recorded in the arXiv licence metadata for this exact version and shown on the abstract page https://arxiv.org/abs/2602.15817v1. Full licence notice: \texttt{licenses/arxiv-2602.15817-CC-BY-4.0.txt}.\par\smallskip

\noindent\textit{Editorial scope.} Counted unit: the article's prop prop (source0.tex lines 1570--1579). The article's complete original proof of it is delivered with it (source0.tex lines 1580--1642, one \texttt{proof} environment of 63 lines). No mathematics is written, repaired or re-derived: the statement and the proof are the article's own lines, in the article's own notation, and no retained line is substituted or re-flowed.  No further source-local context is needed: every cross-reference of every retained line resolves inside the two delivered spans.  Nothing was omitted from any retained span, so the package carries no omission record.  No retained line contains a citation, so the wrapper carries no bibliography and no bibliography entry.  No retained line contains a figure environment, an image-inclusion command or drawing code, so no image is delivered and the package declares no asset dependency.  Editorial formatting only, in addition: an AMS \texttt{amsart} preamble that restates the article's own declarations and macros used by the retained lines, namely the environments \texttt{prop} on separate counters, together with the standard packages those lines need.  The retained environments print in this document as Proposition 1 in order of appearance, whereas the article prints the same items with the numbering of its own full document.  The complete native source of the article as distributed in the arXiv e-print for this exact version is bundled unchanged as \texttt{sources/194831/source0.tex}.
\begin{prop}
    The fixed-point always exists, and is given by 
    \begin{equation}
        q^{(\infty)}(\feas{}=0 | \theta) = \max \left( 0, p^\pi(\feas{}=0 | \theta) - \alpha \frac{ p_{\MeasuredCISet}(\theta) p^\pi(\feas{}=0) }{ p(\theta) } \right),
    \end{equation}  
    where
    \begin{equation}
        p^\pi(\feas{}=0) = \int p^\pi(\feas{}=0 | \theta) p(\theta) d\theta = \E_{\theta \sim p(\theta)} \left[ p^\pi(\feas{}=0 | \theta) \right].
    \end{equation}
\end{prop}

\begin{proof}
Note that $\frac{ (1-\alpha) \rho^{(k)}(\theta) }{ \alpha p_{\MeasuredCISet}(\theta) + (1-\alpha) \rho^{(k)}(\theta) } \in (0, 1)$, so the error rate of the classifier is always upper-bounded by the policy-induced error rate:
\begin{equation}
    q^{(k+1)}(\feas{}=0 | \theta) \leq p^\pi(\feas{}=0 | \theta).
\end{equation}
Moreover, note that $0$ is absorbing for the iterates: if $q^{(k)}(\feas{}=0 | \theta) = 0$, then $\rho^{(k)}(\theta) = 0$ and thus $q^{(k+1)}(\feas{}=0 | \theta) = 0$.

For $\theta$ such that $p^\pi(\feas{}=0 | \theta) > 0$, we can analyze the fixed points of the iterates. The fixed point $q^{(\infty)}$ satisfies
\begin{equation}
    q^{(\infty)}(\feas{}=0 | \theta) = p^\pi(\feas{}=0 | \theta) \frac{ (1-\alpha) \frac{q^{(\infty)}(\feas{}=0|\theta) p(\theta) }{q^{(\infty)}(\feas{}=0)} }{ \alpha p_{\MeasuredCISet}(\theta) + (1-\alpha) \frac{q^{(\infty)}(\feas{}=0|\theta) p(\theta)}{q^{(\infty)}(\feas{}=0)} }.
\end{equation}
Or,
\begin{equation}
    q^{(\infty)}(\feas{}=0 | \theta)
    \left( \alpha p_{\MeasuredCISet}(\theta) + (1 - \alpha) \frac{ q^{(\infty)}(\feas{}=0|\theta) p(\theta) }{ q^{(\infty)}(\feas{}=0) } - (1-\alpha) p^\pi(\feas{}=0|\theta) \frac{p(\theta)}{q^{(\infty)}(\feas{}=0)} \right) = 0
\end{equation}
In other words, at the fixed point, either $q^{(\infty)}(\feas{}=0 | \theta) = 0$ or
\begin{equation}
    \alpha p_{\MeasuredCISet}(\theta) + (1 - \alpha) \frac{ q^{(\infty)}(\feas{}=0|\theta) p(\theta) }{ q^{(\infty)}(\feas{}=0) } - (1-\alpha) p^\pi(\feas{}=0|\theta) \frac{p(\theta)}{q^{(\infty)}(\feas{}=0)} = 0.
\end{equation}
Assume the latter. Then,
\begin{align}
    q^{(\infty)}(\feas{}=0 | \theta)
    &= p^\pi(\feas{}=0 | \theta) \frac{1}{1 + \frac{\alpha}{1-\alpha} \frac{ p_{\MeasuredCISet}(\theta) }{ \rho^{(\infty)}(\theta) } }, \\
    &= p^\pi(\feas{}=0 | \theta) \frac{1}{1 + \frac{\alpha}{1-\alpha} \frac{ p_{\MeasuredCISet}(\theta) q^{(\infty)}(\feas{} = 0) }{ q^{(\infty)}(\feas{}=0 | \theta) p(\theta) } }
\end{align}
Rearranging,
\begin{align}
    q^{(\infty)}(\feas{}=0 | \theta) \left(1 + \frac{\alpha}{1-\alpha} \frac{ p_{\MeasuredCISet}(\theta) q^{(\infty)}(\feas{} = 0) }{ q^{(\infty)}(\feas{}=0 | \theta) p(\theta) } \right)
    &= p^\pi(\feas{}=0 | \theta), \\
    \iff
    q^{(\infty)}(\feas{}=0 | \theta) + \frac{\alpha}{1-\alpha} \frac{ p_{\MeasuredCISet}(\theta) q^{(\infty)}(\feas{} = 0) }{ p(\theta) }
    &= p^\pi(\feas{}=0 | \theta), \\
    \iff
    q^{(\infty)}(\feas{}=0 | \theta)
    &= p^\pi(\feas{}=0 | \theta) - \frac{\alpha}{1-\alpha} \frac{ p_{\MeasuredCISet}(\theta) q^{(\infty)}(\feas{} = 0) }{ p(\theta) }.
\end{align}
Substituting $q^{(\infty)}(\feas{}=0 | \theta)$ into $q^{(\infty)}(\feas{} = 0)$, we get
\begin{align}
    q^{(\infty)}(\feas{} = 0)
    &= \int q^{(\infty)}(\feas{}=0 | \theta) p(\theta) d\theta, \\
    &= \int \left( p^\pi(\feas{}=0 | \theta) - \frac{\alpha}{1-\alpha} \frac{ p_{\MeasuredCISet}(\theta) q^{(\infty)}(\feas{} = 0) }{ p(\theta) } \right) p(\theta) d\theta, \\
    &= \int p^\pi(\feas{}=0 | \theta) p(\theta) d\theta - \frac{\alpha}{1-\alpha} q^{(\infty)}(\feas{} = 0) \int p_{\MeasuredCISet}(\theta) d\theta, \\
    &= p^\pi(\feas{}=0) - \frac{\alpha}{1-\alpha} q^{(\infty)}(\feas{} = 0).
\end{align}
Hence,
\begin{align}
    q^{(\infty)}(\feas{} = 0)
    &= p^\pi(\feas{}=0) - \frac{\alpha}{1-\alpha} q^{(\infty)}(\feas{} = 0), \\
    \implies
    q^{(\infty)}(\feas{} = 0)
    &= (1-\alpha) p^\pi(\feas{}=0).
\end{align}
Substituting back, we get the fixed point
\begin{equation}
    q^{(\infty)}(\feas{}=0 | \theta)
    = p^\pi(\feas{}=0 | \theta) - \alpha \frac{ p_{\MeasuredCISet}(\theta) p^\pi(\feas{}=0) }{ p(\theta) }.
\end{equation}
Since $q^{(\infty)}(\feas{}=0 | \theta) \geq 0$, we have that
\begin{equation}
    q^{(\infty)}(\feas{}=0 | \theta) = \max \left( 0, p^\pi(\feas{}=0 | \theta) - \alpha \frac{ p_{\MeasuredCISet}(\theta) p^\pi(\feas{}=0) }{ p(\theta) } \right).
\end{equation}
\end{proof}

\end{document}
